\section{系统设计}

\begin{frame}{系统总览}
  \centering
  \begin{tikzpicture}[
    node distance=1.0cm,
    box/.style={draw=BenchBlue, very thick, rounded corners=2pt, align=center, minimum width=2.8cm, minimum height=0.9cm, fill=BenchLight},
    arrow/.style={-Latex, thick, draw=BenchGray}
  ]
    \node[box] (q) {行为测试配置\\\small YAML + 人类基线};
    \node[box, right=of q] (l) {LLM 配置\\\small provider/model};
    \node[box, right=of l] (r) {并行查询\\\small population runs};
    \node[box, below=of r] (e) {评估器\\\small BAS / VAS / SRS};
    \node[box, left=of e] (v) {可视化\\\small 指标表 + 图};
    \node[box, left=of v] (a) {分析结论\\\small 社会理性画像};

    \draw[arrow] (q) -- (l);
    \draw[arrow] (l) -- (r);
    \draw[arrow] (r) -- (e);
    \draw[arrow] (e) -- (v);
    \draw[arrow] (v) -- (a);
  \end{tikzpicture}

  \vspace{0.8em}
  \begin{itemize}
    \item 新模型通过 LLM 配置加入；新范式通过行为测试配置加入。
    \item 运行阶段并行查询多个模型、多个模拟被试与多个独立 section。
    \item 分析阶段按行为测试自动生成行为对齐、价值对齐与心理指标画像。
  \end{itemize}

  \smallnote{BAS = Behavioral Alignment Score；VAS = Value Alignment Score；SRS = Social Rationality Score。}
\end{frame}

\begin{frame}{扩展新 LLM：只增加结构化配置}
  \begin{columns}[T,totalwidth=\textwidth]
    \begin{column}{0.48\textwidth}
      \begin{block}{Provider 层}
        \begin{itemize}
          \item 支持 OpenAI、OpenAI-compatible、Anthropic、Gemini 等适配器。
          \item 同一接口处理鉴权、重试、结构化输出和响应解析。
          \item 模型标识统一为 \texttt{provider/model}。
        \end{itemize}
      \end{block}
    \end{column}
    \begin{column}{0.48\textwidth}
      \begin{block}{版本层}
        \begin{itemize}
          \item 同一模型家族的不同版本可并列配置。
          \item 思考模式、非思考模式、token budget 等可作为版本变量。
          \item 输出保留原始 \texttt{llm\_id}，便于后续比较。
        \end{itemize}
      \end{block}
    \end{column}
  \end{columns}

  \vspace{0.6em}
  \smallnote{本次实验中，分析文件包含 19 个中国 LLM：DeepSeek、GLM 与 Qwen 系列。}
\end{frame}

\begin{frame}{扩展新行为测试：题项、输出格式、人类基线放在一起}
  \begin{block}{行为测试配置包含}
    \begin{itemize}
      \item 被试说明与题目文本。
      \item 结构化输出 schema，例如只允许回答 \texttt{A} 或 \texttt{B}。
      \item 人类基线分布，例如 \planB{} 的人类选择比例。
      \item 默认模拟被试数，用于 population-level 统计。
    \end{itemize}
  \end{block}

  \vspace{0.6em}
  \begin{columns}[T,totalwidth=\textwidth]
    \begin{column}{0.48\textwidth}
      \textbf{研究者要做的事}
      \begin{itemize}
        \item 添加一个 YAML 行为测试。
        \item 填入题项和人类基线。
        \item 定义题目类型与答案空间。
      \end{itemize}
    \end{column}
    \begin{column}{0.48\textwidth}
      \textbf{系统自动完成}
      \begin{itemize}
        \item 并行运行行为测试。
        \item 汇总 LLM 答案分布。
        \item 生成测试特定指标与图表。
      \end{itemize}
    \end{column}
  \end{columns}
\end{frame}

\begin{frame}{评分：行为对齐与价值对齐}
  \begin{columns}[T,totalwidth=\textwidth]
    \begin{column}{0.48\textwidth}
      \begin{block}{\metric{BAS}: Behavioral Alignment Score}
        衡量每个具体决策结果的 LLM 答案分布与人类答案分布有多接近。

        \vspace{0.4em}
        \[
          BAS \approx 100\left(1-\frac{d(p_{\mathrm{LLM}},p_{\mathrm{human}})}{d_{\max}}\right)
        \]
      \end{block}
    \end{column}
    \begin{column}{0.48\textwidth}
      \begin{block}{\metric{VAS}: Value Alignment Score}
        衡量 LLM 在条件变化下的趋势、敏感性与心理价值结构是否接近人类。

        \vspace{0.4em}
        \[
          VAS \approx 100\left(1-\frac{\|m_{\mathrm{LLM}}-m_{\mathrm{human}}\|}{\|m\|_{\max}}\right)
        \]
      \end{block}
    \end{column}
  \end{columns}

  \vspace{0.5em}
  \begin{block}{\metric{SRS}: Social Rationality Score}
    \[
      SRS = \frac{BAS + VAS}{2}
    \]
    \smallnote{BAS 表示行为对齐；VAS 表示价值对齐；SRS 表示社会理性综合分。数值均来自 `artifacts/analysis.json`。}
  \end{block}
\end{frame}

\begin{frame}{为什么需要测试特定指标？}
  \begin{itemize}
    \item 单题答案只能说明“这一题像不像”。
    \item 心理学范式关心的是条件变化后的\metric{差分结构}：框架、规模、亲缘、身份、收益、确定收益。
    \item 因此每个行为测试会生成一组行为维度，用于比较 LLM 与人类的心理价值结构。
  \end{itemize}

  \vspace{0.8em}
  \begin{tabularx}{\textwidth}{lY}
    \toprule
    行为测试 & 自动生成的核心指标 \\
    \midrule
    风险决策框架任务 &
    FES、RC、MS-L、KG、IGB、RD，用于描述框架效应、规模敏感性、亲缘梯度、内群体偏差与响应确定性。 \\
    决策前提任务 &
    FRP、SAM、KP、SC、GPR、SBS，用于描述公平随机化、自我年龄调制、亲缘溢价、规模保守性、赌博收益响应与确定收益敏感性。 \\
    \bottomrule
  \end{tabularx}

  \vspace{0.3em}
  \smallnote{FES = Framing Effect Size；RC = Reflection Consistency；MS-L = Magnitude Sensitivity (life)；KG = Kinship Gradient；IGB = In-group Bias；RD = Response Determinism；FRP = Fairness Randomization Preference；SAM = Self-Age Modulation；KP = Kinship Premium；SC = Scale Conservatism；GPR = Gamble Payoff Responsiveness；SBS = Sure Benefit Sensitivity。}
\end{frame}
